\documentclass[10pt,a4paper]{amsart}
\usepackage[margin=19mm]{geometry}
\usepackage{amsmath,amssymb,mathtools}
\usepackage[scr=rsfs,cal=euler]{mathalfa}
\usepackage{xfrac}
\usepackage[unicode,hidelinks,bookmarks=false]{hyperref}
\newcommand{\ie}{i.e.\ }
\newcommand{\cf}{cf.\ }
\newcommand{\resp}{respectively\ }
\newcommand{\Subsection}{\subsection}
\providecommand{\nobreakdash}{\nobreak\hbox{-}}
\setlength{\emergencystretch}{2em}
\multlinegap0pt
\usepackage{bm}
\usepackage{bbm}
\usepackage{dsfont}
\usepackage{xspace}
\usepackage{cancel}
\usepackage{mathtools}
\usepackage{amsthm}
\makeatletter
\@ifundefined{Definition}{\newtheorem{Definition}{Definition}}{}
\@ifundefined{Proposition}{\newtheorem{Proposition}[Definition]{Proposition}}{}
\makeatother
\newcommand{\R}{\mathbb{R}}
\newcommand{\T}{\mathbb{T}}
\newcommand{\al}{\alpha}
\newcommand{\eps}{\varepsilon}
\title[On the Dynamics of Invariant Graphs for Dissipative Twist Ma]{On the Dynamics of Invariant Graphs for Dissipative Twist Maps}
\author{Qi Li, Lin Wang}
\date{Source: arXiv:2506.04938v2 (CC BY 4.0)}
\begin{document}
\maketitle

\noindent\emph{Source and licence.} On the Dynamics of Invariant Graphs for Dissipative Twist Maps. Version arXiv:2506.04938v2 (CC BY 4.0), distributed under the Creative Commons Attribution 4.0 International licence, as recorded in the arXiv licence metadata for this exact version and shown on the abstract page https://arxiv.org/abs/2506.04938v2. Full rights notice: licenses/arxiv-2506.04938-CC-BY-4.0.txt.\par\smallskip

\noindent\textit{Editorial scope.} Counted unit: the Proposition whose source label is \texttt{lipo} in the article (source0.tex lines 605--608, source label \texttt{lipo}), delivered together with the article's own complete original proof of it (source0.tex lines 609--682, one \texttt{proof} environment of 74 lines). No mathematics is written, repaired or re-derived: statement and proof are the article's own text line for line. Required context retained uncounted: hf1 (lines 565--570); nhim (lines 585--590). Every definition, display and label of those lines is retained unchanged. Omitted from this package and not redistributed: 1--564, 571--584, 591--604, 683--2092. Nothing inside a retained excerpt is omitted or altered: every retained body is delivered exactly as the article's lines are written, so no omission receipt of any kind accompanies this package. Citations: the retained excerpts carry no citation command at all, so no bibliography entry of the article is transcribed and no reference is redistributed. Reference adaptation: none was needed and none was made; every reference inside the delivered unit resolves inside it, so no unresolved reference or citation is produced. Printed slips kept literal: no printed word, symbol, hypothesis or number of the counted statement or of its proof was altered. Layout and class: the article's own class is \texttt{amsart} with packages including amsfonts, amssymb, amsmath, amscd, euscript, verbatim, array, geometry, anysize, fancyhdr, indentfirst, graphicx, color, ifpdf; the wrapper is the lane's pinned \texttt{amsart} editorial wrapper at 10pt on A4, which restates the theorem-like environments the retained text needs and adds the article's own macro definitions that the retained text needs (\texttt{\textbackslash R}, \texttt{\textbackslash T}, \texttt{\textbackslash al}, \texttt{\textbackslash eps}), with extra wrapper packages (bm, bbm, dsfont, xspace, cancel, mathtools). The delivered numbering is the unit's own. Assets: no retained span contains a figure environment, a graphics-inclusion command, a PDF-page inclusion, a table or a TikZ picture; no image file is copied into this package, the package declares no asset dependency and no image file is redistributed with it. Licence: version arXiv:2506.04938v2 of this article is distributed under the Creative Commons Attribution 4.0 International licence, as recorded in the arXiv licence metadata for this exact version and shown on the abstract page https://arxiv.org/abs/2506.04938v2. A full rights notice is delivered at licenses/arxiv-2506.04938-CC-BY-4.0.txt. Characters: every retained excerpt is pure ASCII, so the delivered engine, XeLaTeX with its default Latin Modern text font, reports no missing character.
\begin{Proposition}\label{hf1}
The map $F^\phi_{\al}$ admits a $C^0$-invariant graph $\tilde{\Gamma} := \{(x, \tilde{\Psi}_{\alpha}(x)) \mid x \in \mathbb{R}\}$ if and only if there exists $g_{\alpha} \in D^r(\mathbb{T})$ satisfying the functional equation
\begin{equation}\tag{A}\label{eq:main}
\frac{1}{1+\lambda}g_{\alpha}(x) + \frac{\lambda}{1+\lambda}g_{\alpha}^{-1}(x) = x + \frac{1}{1+\lambda}\left((1-\lambda)\alpha_1 + \lambda\alpha_2 + \phi(x)\right) \quad \forall x \in \mathbb{R}.
\end{equation}
\end{Proposition}

\begin{Proposition}\label{nhim}
For any $r \in [1,+\infty)$ and $0 < \eps \ll 1$, there exists $\delta = O(\eps)$ such that if $\|\phi\|_{C^r} \leq \delta$, then the perturbed map $F^{\phi}_\alpha$ admits a $C^r$-invariant graph $\tilde{\Gamma}_\alpha^\phi := \{(x, \tilde{\Psi}_{\alpha}(x)) \mid x \in \mathbb{R}\}$ with the estimate
\[
\left\|\tilde{\Psi}_\alpha - \frac{\alpha_2}{1-\lambda}\right\|_{C^r} \leq \eps.
\]
\end{Proposition}

\begin{Proposition}\label{lipo}
 If $\|\phi\|_{C^1}\leq \delta_0$ where $\delta$ is the constant determined by Proposition \ref{nhim}, then the function $\tilde{\Psi}_\alpha(x)$ is uniformly Lipschitz with respect to the parameter $\alpha:=(\alpha_1,\alpha_2)\in \R^2$. More precisely, for each $x\in \R$ and $\al,{\al'}\in \R$,
\[|\tilde{\Psi}_\al(x)-\tilde{\Psi}_{\al'}(x)|\leq \frac{1}{1-\lambda}|\al_2-\al'_2|.\]
\end{Proposition}

\begin{proof}
Fix the perturbation $\phi:\T\to\R$. For simplicity, we denote $F_\alpha := F^{\phi}_\alpha$. Given $\alpha' := (\alpha'_1, \alpha'_2) \in \R^2$, we consider the transformation $\Phi:\R^2\to\R^2$ defined by
\[
\Phi(x,y) = (x, y - \Psi_{\alpha'}(x)).
\]
Its inverse is given by
\[
\Phi^{-1}(x,y) = (x, y + \Psi_{\alpha'}(x)).
\]
Define the conjugated map
\[
K_{\alpha}(x,y) := \Phi \circ F_\alpha \circ \Phi^{-1}(x,y).
\]
A direct calculation yields
\begin{align*}
K_{\alpha}(x,y) = \biggl(x +  \alpha_1+\lambda(y + \Psi_{\alpha'}(x)) +\phi(x), \lambda(y + \Psi_{\alpha'}(x)) + \alpha_2 + \phi(x) - \Psi_{\alpha'}(x)\biggr).
\end{align*}

Let $\chi(x,y,\alpha) := \pi_2 K_\alpha(x,y)$. Since $\Psi_{\alpha'}$ is an invariant graph for $F_{\alpha'}$, Proposition~\ref{hf1} implies
\[
F_{\alpha'}(x, \Psi_{\alpha'}(x)) = (g_{\alpha'}(x), \Psi_{\alpha'}(g_{\alpha'}(x))),
\]
and consequently $\chi(x,0,\alpha') = 0$ for all $x \in \R$. Define
\[
G(x,y,\alpha) := \chi(x,y,\alpha) - y.
\]
The partial derivatives satisfy:
\begin{align}
\frac{\partial \chi}{\partial y}(x,y,\alpha) = \lambda,\quad \frac{\partial \chi}{\partial \alpha_1}(x,y,\alpha) = 0,\quad \frac{\partial \chi}{\partial \alpha_2}(x,y,\alpha) = 1.
\end{align}

It follows that
\begin{equation}\label{contrpo}
0 < \frac{\partial \chi}{\partial y}(x,y,\alpha)  < 1
\end{equation}
for all $x \in \R$.

Applying the implicit function theorem, there exists a neighborhood $U(\alpha')$ and a $C^1$ function $\sigma:\R \times U(\alpha') \to \R$ such that
\[
\chi(x, \sigma(x, \alpha), \alpha) = \sigma(x, \alpha).
\]
This function satisfies:
\begin{equation}
 \sigma(x, \alpha') = 0,\quad
\left|\frac{\partial \sigma}{\partial \alpha_1}\right| = \left|\frac{\frac{\partial \chi}{\partial \alpha_1}}{1 - \frac{\partial \chi}{\partial y}}\right|=0,\quad
\frac{\partial \sigma}{\partial \alpha_2}= \left|\frac{\frac{\partial \chi}{\partial \alpha_2}}{1 - \frac{\partial \chi}{\partial y}}\right|=\frac{1}{1-\lambda}.
\end{equation}

Consequently, for each $\alpha \in U(\alpha')$ and all $x \in \R$, we have the estimate
\[
|\sigma(x, \alpha)| \leq \frac{1}{1-\lambda}|\alpha_2 - \alpha'_2|.
\]


From \eqref{contrpo}, we obtain the non-expensive property:
\begin{align*}
|\pi_2 K_\alpha(x,y) - \pi_2 K_\alpha(x, \sigma(x, \alpha))|
&= |\chi(x,y,\alpha) - \chi(x, \sigma(x, \alpha), \alpha)| \\
&\leq \left|\frac{\partial \chi}{\partial y}\right| |y - \sigma(x, \alpha)| \\
&\leq |y - \sigma(x, \alpha)|.
\end{align*}

This shows that the set
\[
\Omega_1 := \{(x,y) \mid |y| \leq \frac{1}{1-\lambda}|\alpha_2 - \alpha'_2|\}
\]
is an attractor block for $K_\alpha$. Through conjugation, the corresponding set
\[
\Omega_2 := \{(x,y) \mid |y - \Psi_{\alpha'}(x)| \leq \frac{1}{1-\lambda}|\alpha_2 - \alpha'_2|\}
\]
is an attractor block for $F_\alpha$. Since the graph of $\Psi_\alpha$ is an attractor for $F_\alpha$, it must be contained in $\Omega_2$.

This completes the proof of Proposition~\ref{lipo}.
\end{proof}

\end{document}
